\documentclass[11pt]{article}

\usepackage[margin=1in]{geometry}
\usepackage{amsmath,amssymb,amsthm}
\usepackage{hyperref}
\usepackage{xurl}
\usepackage{microtype}
\usepackage{enumitem}
\usepackage{booktabs}
\usepackage{listings}
\usepackage{float}
\usepackage[nameinlink,noabbrev]{cleveref}

\theoremstyle{definition}
\newtheorem{definition}{Definition}
\theoremstyle{remark}
\newtheorem{remark}{Remark}

\title{Closed-View Auditing and CPPCs (Addendum):\\Multi-Stream Monitoring Notes and Schema Excerpts}
\author{}
\date{June 18, 2026 (rev0898 artifact-boundary repair)}

\begin{document}
\maketitle

\begin{abstract}
This companion note collects only the supplementary control-plane packet for the main closed-view auditing paper (State~3):
(i) a compact many-detectors monitoring note using e-values and scan-inflation control,
(ii) a short CPPC schema excerpt,
(iii) a reference-linter excerpt, and
(iv) two reference CPPC card fragments.
It is intentionally narrow: we do not restate the closed-view monitoring theorem/root claim, evaluator notarization surface, runtime-plan conformance receipts, or public release-lineage packaging.
The note exists so State~3 can stay referee-tight while still pointing to one small, auditable appendix-level packet.
\end{abstract}

\paragraph{Scope.}
Addendum to State~3.\cite{series:state3}

\paragraph{Imports.}
We import the CPPC concept, the measurement-firewall pattern, and the alert-oracle framing from State~3.\cite{series:state3}
The surrounding evaluation / publication stack is imported rather than re-owned here: evaluator notarization belongs to Eval~1, runtime-plan conformance receipts belong to Operational~B, and public claim packaging / release lineage belong to ABOM/OINL plus Release~A.\cite{series:evalnotary,series:opB,series:abomoinl,series:releaseproperty}
Only the inline artifacts printed in this note are claimed; no request-only support bundle is part of the evidence.

\paragraph{Exports.}
A compact supplementary packet: one many-detectors monitoring note, one short CPPC schema excerpt, one inline reference-linter excerpt, and two reference card fragments.

\paragraph{Boundary.}
This addendum owns only those supplementary monitoring and schema artifacts.
It does \emph{not} own the closed-view monitoring contract and alert-oracle tax (State~3), evaluator notarization / anti-grinding evidence release (Eval~1), runtime-plan pinning / PTL / epoch receipts (Operational~B), or public claim packaging / release lineage (ABOM/OINL plus Release~A).\cite{series:state3,series:evalnotary,series:opB,series:abomoinl,series:releaseproperty}

\paragraph{Exported object.}
The exported object is a small support packet: one many-detector scan-inflation recipe that a referee can quote, one schema/linter excerpt that makes ``machine-checkable CPPC'' concrete, and two short reference card fragments showing how the manifest vocabulary lands on real systems without implying that a larger schema or tooling bundle exists.

\paragraph{Earliest sufficient stop.}
This addendum is complete once a reader can cite the many-detector recipe or the CPPC schema/card excerpt as explanatory support for a State~3 claim.
If a deployment needs the underlying closed-view theorem packet, a notarized evaluation log, a runtime conformance receipt, or a public release-lineage bundle, the reader should leave this addendum and use State~3, Eval~1, Operational~B, and ABOM/OINL plus Release~A directly.\cite{series:state3,series:evalnotary,series:opB,series:abomoinl,series:releaseproperty}

\paragraph{Later terse-paper citation posture.}
For later terse papers under space pressure, default to State~3 for the closed-view monitoring contract and public CPPC claim; cite this addendum only when the exact support packet exported here --- the many-detector monitoring note, the schema excerpt, the inline linter rule, or the reference card fragments --- is itself live.
If the live question is only the public monitoring theorem or alert-oracle claim, stop at State~3.
If the live question is instead ``where is the concrete schema or many-detector recipe that supports that claim?'', widen to this addendum.
If the question becomes evaluator notarization, runtime conformance, or public release-lineage comparison, leave this addendum and reopen the neighboring owners.
The maintained worked example and paired cutover rule are the general model for that stop / widen / reopen discipline.\cite{series:workedexample,series:pairedterminalcutoverrules}

\paragraph{Role in the branched state-family route.}
In the state family this addendum is never the theorem root, never the anonymous-DHT committee-contact specialization, and never the public monitoring card itself.  It is only the optional final support packet beneath State~3.  A later terse paper may therefore reach this note from more than one lower branch --- a generic State~1 packet, a State~1A calibration sentence, or a State~2 MUCC sentence may all sit underneath the same monitored deployment slice --- but the widening rule is the same in every case: stop at State~3 for the public monitoring claim, and widen here only if the concrete schema/linter/many-detector support packet must itself be quoted.  This addendum explains a monitoring card that sits above a branchy family; it never replaces any lower branch or the main monitoring stop.\cite{series:state1,series:state1addendum,series:state2,series:state3}

\paragraph{Neighboring-layer stop map.}
Table~\ref{tab:cppc-addendum-stopmap} states the smallest handoff that keeps this addendum narrow.
A referee can stay here for the many-detector recipe, schema excerpt, inline linter rule, and reference card fragments, but should widen immediately once the question becomes one of theorem justification, notarized evidence, runtime-plan conformance, or public release packaging.

\begin{table}[H]
\centering
\begin{tabular}{@{}p{0.34\linewidth}p{0.58\linewidth}@{}}
\toprule
Question class & Leave this addendum for \\
\midrule
What closed-view monitoring contract or alert-oracle tax justifies the claimed witness discipline? & State~3, which owns the closed-view monitoring theorem/root claim and the alert-oracle pricing rule. \\
Which evaluator log fixes the detector set, evidence window, and anti-grinding attempt record? & Eval~1, which owns notarized evaluator manifests, gap-free attempt prefixes, and verifier reports. \\
Did the monitored deployment actually follow the pinned runtime plan and publication cadence? & Operational~B, which owns plan pinning, PTL inclusion, and epoch-conformance receipts. \\
Which public card / label / release receipt makes the monitoring claim publishable across versions? & ABOM/OINL plus Release~A, which own public claim packaging and release-lineage comparison. \\
\bottomrule
\end{tabular}
\caption{Neighboring-layer stop map for the State~3 CPPC addendum. This note owns only the supplementary control-plane packet.}
\label{tab:cppc-addendum-stopmap}
\end{table}

\paragraph{A minimal public CPPC support-packet card.}
When a later short paper only needs to answer \emph{which supplementary control-plane packet made this State~3 monitoring claim concrete?}, it should stop at the small card below plus the tiny support-packet lookup drill that follows.

\begin{table}[H]
\centering
\begin{tabular}{@{}p{0.30\linewidth}p{0.60\linewidth}@{}}
\toprule
Support-packet field & Minimal public content a later terse paper may quote here \\
\midrule
Packet claim & This addendum exports only the State~3 support packet: one many-detector monitoring note, one CPPC schema excerpt, one inline linter rule, and two reference card fragments. \\
Many-detector anchor & The packet's monitoring note is the e-process / e-BH scan-inflation recipe in \Cref{app:multi}; it explains how multiple closed-view detectors are monitored without turning this addendum into a second theorem paper. \\
Schema excerpt anchor & The public schema excerpt in \Cref{app:cppc} exposes the CPPC fields a later paper may name directly: \texttt{witness\_channels\_present}, \texttt{cadence}, \texttt{view\_accountant}, \texttt{budget\_beacon}, \texttt{transparency}, and \texttt{compliance}. \\
Linter anchor & The inline reference-linter excerpt below shows the machine-checkable support rule directly: required keys present, enum check on witness channels, and artifact-existence / hash checks when provided. \\
Reference card anchors & The two public card fragments are the Nym-style and Tor-style reference excerpts below; they are explanatory support for the CPPC vocabulary, not new theorem or release objects. \\
Wrapper-refresh cutover & Keep this same card only if a successor wrapper preserves the same many-detector recipe anchor, the same schema/linter excerpt, and the same reference card-fragment role. \\
How to cite tersely & Quote this card only when the live issue is the concrete support packet itself; otherwise stop at State~3 for the public monitoring claim. \\
Do not widen silently & Evaluator notarization, runtime-plan conformance, and public release lineage still belong to Eval~1, Operational~B, and ABOM/OINL plus Release~A. \\
\bottomrule
\end{tabular}
\caption{A minimal public CPPC support-packet card. This is the smallest public answer later terse papers can quote directly when the live issue is the concrete many-detector note, schema excerpt, inline linter rule, or reference CPPC card fragments rather than the main State~3 monitoring contract.}
\label{tab:cppc-addendum-card}
\end{table}

\paragraph{Tiny support-packet lookup drill.}
Suppose a later short paper only needs to justify the sentence ``the monitored deployment published a CPPC support packet with \texttt{budget\_beacon} and \texttt{auditing\_alerts} channels, an RDP odometer, and a machine-checkable linter path.''
The smallest public support object is exactly this addendum card plus the schema excerpt in \Cref{app:cppc}, the inline reference-linter rule below, and the Nym-style reference fragment below.
That packet already names \texttt{schema\_version = 1.1}, \texttt{witness\_channels\_present = [budget\_beacon, auditing\_alerts, timing\_and\_metadata]}, a one-minute epoch / daily delayed publication cadence, an RDP odometer budget block, and the rule that required keys, enum membership, and declared artifact existence/hash checks must pass when provided.
A later terse paper can therefore cite one concrete CPPC support packet without reopening State~3's theorem root, Eval~1's notarized evidence packet, Operational~B's runtime receipts, or Release~A's lineage object.

\paragraph{Tiny machine-checkable support sentence.}
The smallest honest sentence a later terse paper may now quote directly is: \emph{``The State~3 support packet fixes a CPPC \texttt{schema\_version = 1.1}, requires \texttt{witness\_channels\_present}, \texttt{cadence}, \texttt{view\_accountant}, \texttt{budget\_beacon}, \texttt{transparency}, and \texttt{compliance}, and checks enum membership plus declared artifact existence/hash when provided.''}
That sentence is narrower than the main monitoring claim but more auditable than a decorative card fragment because every quoted field now appears in the inline schema/linter support packet itself.

\paragraph{Paired route back to State~3.}
When a later terse paper needs both the public monitoring-contract sentence and the concrete support-packet sentence, the paired route is explicit: State~3 owns the signed CPPC card and the public alert boundary, while this addendum owns only the optional support packet that explains how the many-detector recipe, schema excerpt, inline linter rule, and reference card fragments make that card concrete.  This addendum therefore does not replace the main monitoring contract; it accompanies it only when the exact support packet is itself live.\cite{series:state3}

\paragraph{Tiny paired drill.}
Suppose a later short paper needs to say both that \texttt{cppc.r1} exported weekly beacons, one delayed coarse alert family, and a $6.0$-bit public-alert cap, and that the concrete support packet used the addendum's e-process / e-BH scan-inflation note plus the schema/linter excerpt naming \texttt{witness\_channels\_present}, \texttt{cadence}, \texttt{view\_accountant}, \texttt{budget\_beacon}, \texttt{transparency}, and \texttt{compliance}.
The first sentence belongs to State~3.
The second belongs here.
That split is the smallest honest paired citation: the support packet explains the main card, but it does not replace it.\cite{series:state3}

\paragraph{Same support-packet / refreshed-wrapper cutover.}
A later terse paper may keep this exact CPPC support-packet card across a successor wrapper refresh only when the live packet is unchanged: the many-detector note still points to the same e-process / e-BH scan-inflation recipe, the public schema excerpt still carries the same witness-channel / cadence / accountant / budget-beacon / transparency / compliance fields, the inline linter rule still checks the same required-key / enum / artifact-existence conditions, and the public reference fragments are still playing the same explanatory role.
If the successor wrapper changes the many-detector guidance, changes which CPPC fields or enum values are publicly exported, changes the linter rule that makes the packet machine-checkable, or swaps in different reference fragments to support the public claim, then the shortcut fails closed and the reader should reopen this addendum rather than silently reusing the old card.

\paragraph{Tiny successor support drill.}
Suppose a successor release changes only the ABOM / release-lineage wrapper while preserving the same multi-detector recipe in \Cref{app:multi}, the same schema excerpt fields and witness-channel enum in \Cref{app:cppc}, the same inline linter rule, and the same Nym-style / Tor-style reference fragments.
Then a later terse paper may honestly keep this same support-packet card and say that only the wrapper changed.
But if the successor packet instead adds a new public witness channel, changes the cadence or accountant lines in the exported excerpt, or tightens the linter so different compliance artifacts are now required, the support packet itself has changed and the later paper must reopen this addendum.

\appendix


\section{Multi-stream monitoring note (e-values)}
\label{app:multi}
Stateful anonymity deployments often monitor many candidate witnesses (many committees/keys/relays/metrics).
A concise way to control ``scan inflation'' is to attach each detector an \emph{e-process} (a nonnegative supermartingale with unit expectation under the null) and apply e-value multiple testing.

\paragraph{E-values and e-BH.}
E-values generalize likelihood-ratio betting scores; they compose naturally under optional stopping and enable dependence-robust FDR control via the e-BH procedure~\cite{ramdas2024evalues,wang2022ebh}.
Operationally, this lets an operator run many closed-view detectors while bounding the expected false discovery proportion.

\paragraph{Stopping-time subtlety.}
When alerts are raised at data-dependent global stopping times (``whenever anything looks bad''), additional care is required; the stopped e-BH procedure is designed for this regime~\cite{wang2025sebH}.
The main paper develops the deployment-side implication: \emph{publishing} the alert time/window can itself be a witness stream and should be priced.\cite{series:state3}



\section{CPPC schema (excerpt) and example cards}
\label{app:cppc}
To keep this note referee-tight, we include only a short \emph{excerpt} of the CPPC schema, one inline reference-linter rule, and two reference card fragments.
No fuller JSON schema, complete card corpus, or machine-checking tool is shipped or claimed.  The excerpt below is the entire support artifact owned by this note.

\subsection{Schema excerpt (v1.1)}
\begin{lstlisting}
{
  "schema_version": "1.1",
  "system": "string",
  "unit_of_privacy": { "...": "..." },
  "witness_streams": [ "string", "..." ],
  "witness_channels_present": [
    "semantic_outputs",
    "budget_beacon",
    "transparency_log",
    "timing_and_metadata",
    "auditing_alerts"
  ],
  "cadence": { "epoch": "...", "publication": "..." },
  "view_accountant": {
    "notion": "max-divergence | RDP | f-DP/GDP",
    "composition": "...",
    "budget": { "...": "..." }
  },
  "budget_beacon": { "cadence": "...", "format": "..." },
  "transparency": { "log": "...", "signatures": "..." },
  "compliance": { "mode": "...", "artifacts": [ "..." ] }
}
\end{lstlisting}

\subsection{Reference linter rule (excerpt)}
To keep the support packet audit-ready without falsely claiming a shipped tooling tree, we include a tiny inline reference-linter rule that checks a CPPC JSON for required fields and basic internal consistency.
The intent is to make ``machine-checkable'' concrete without claiming an absent tooling tree or requiring any particular SBOM or transparency-log ecosystem.

\begin{lstlisting}
require_keys(card, [
  "schema_version",
  "witness_channels_present",
  "cadence",
  "view_accountant",
  "budget_beacon",
  "transparency",
  "compliance"
])
assert subset(card["witness_channels_present"], ALLOWED_CHANNEL_ENUM)
for artifact in card["compliance"].get("artifacts", []):
    assert exists_under(ROOT, artifact)
    if has_declared_sha256(card, artifact):
        assert sha256_matches(ROOT, artifact, declared_sha256(card, artifact))
warn_if_missing(card["transparency"], "log")
warn_if_missing(card["transparency"], "signatures")
\end{lstlisting}

In the stop map above, that inline rule remains only a support artifact: it makes the CPPC vocabulary auditable, but it does not by itself settle whether the monitored deployment satisfied the State~3 theorem assumptions, matched the pinned runtime plan, or was packaged into a publishable release claim.



\subsection{Example CPPC (Nym-style, reference excerpt)}
\begin{lstlisting}
{
  "schema_version": "1.1",
  "system": "nym-mixnet (reference fragment)",
  "unit_of_privacy": { "message": "1 packet-equivalent" },
  "witness_channels_present": [
    "budget_beacon",
    "auditing_alerts",
    "timing_and_metadata"
  ],
  "cadence": {
    "epoch": "1 min",
    "publication": "daily aggregates; delayed by 7d"
  },
  "witness_streams": [
    "mix-queue depth (bucketed)",
    "node reputation updates (coarsened)",
    "drift alerts (delayed, bucketed)"
  ],
  "view_accountant": {
    "notion": "RDP",
    "composition": "odometer (adaptive cadence)",
    "budget": {
      "epsilon_rdp": 3.0,
      "delta": 1e-6,
      "horizon": "30d rolling"
    }
  },
  "budget_beacon": {
    "cadence": "daily",
    "payload": "current spend + remaining budget (coarsened)"
  },
  "transparency": {
    "log": "append-only merkle log",
    "signatures": "operator key + optional attestation"
  },
  "compliance": {
    "mode": "attestation",
    "artifacts": [ "cppc.json", "build hash", "release hashes" ]
  }
}
\end{lstlisting}



\subsection{Example CPPC (Tor-style, reference excerpt)}
\begin{lstlisting}
{
  "schema_version": "1.1",
  "system": "tor-relay-selection (reference fragment)",
  "unit_of_privacy": { "client": "one long-lived guard set" },
  "witness_channels_present": [
    "semantic_outputs",
    "budget_beacon",
    "transparency_log"
  ],
  "cadence": {
    "epoch": "1 consensus period",
    "publication": "public consensus + coarse metrics"
  },
  "witness_streams": [
    "selection weights (public consensus)",
    "bandwidth measurement summaries (coarse)",
    "guard selection randomness"
  ],
  "view_accountant": {
    "notion": "max-divergence",
    "composition": "filter (abort-safe)",
    "budget": { "epsilon": 2.0, "delta": 1e-6, "horizon": "90d" }
  },
  "budget_beacon": {
    "cadence": "each consensus period",
    "payload": "monotone spend counter + checkpoint hash"
  },
  "transparency": {
    "log": "public merkle log (checkpointed)",
    "signatures": "authority keys"
  },
  "compliance": {
    "mode": "hybrid",
    "artifacts": [
      "cppc.json",
      "consensus hashes",
      "declared release hashes"
    ]
  }
}
\end{lstlisting}







\begin{thebibliography}{99}

\bibitem{series:state3}
Anonymous.
\newblock Closed-View Auditing for Stateful Anonymity: Measurement Firewalls and CPPCs.
\newblock AnonDHT State series Paper~3, 2026. (This archive.)

\bibitem{series:contractobjects}
Anonymous.
\newblock Contract Objects for Anonymous-DHT Privacy Budgets and Claims.
\newblock Series synthesis note (Synthesis~0), 2026. (This archive.)

\bibitem{series:abomoinl}
Anonymous.
\newblock ABOM/OINL: Audit BOM and Outcome Identification Noise Label for Anonymous DHT Deployments.
\newblock Anonymity series Paper~C, 2026. (This archive.)

\bibitem{series:evalnotary}
Anonymous.
\newblock Notarized Anonymity Certificates Under Retries.
\newblock Evaluation series Paper~1, 2026. (This archive.)

\bibitem{series:opB}
Anonymous.
\newblock Auditable Trace Privacy for Anonymous DHTs: Runtime Conformance, Plan Pinning, and Epoch Receipts.
\newblock Operational series Paper~B, 2026. (This archive.)

\bibitem{series:releaseproperty}
Anonymous.
\newblock Release Property Ledger Receipts for Anonymity Claims.
\newblock Release \& destination Paper~A, 2026. (This archive.)

\bibitem{series:workedexample}
Anonymous.
\newblock Worked Example: Instantiating a Tiered Reader-Privacy Receipt.
\newblock Synthesis~17, The-One-True-Archive (2026).

\bibitem{series:pairedterminalcutoverrules}
Anonymous.
\newblock Paired Terminal Cutover Rules for Stable Review Spines: When Later Terse Papers Must Stop, Widen, or Reopen.
\newblock Series synthesis note (Synthesis~75), 2026. (This archive.)

\bibitem{series:artifactpolicy}
Anonymous.
\newblock Artifact policy and supporting materials for the anonymity/AnonDHT series.
\newblock 2026.

\bibitem{ramdas2024evalues}
Aaditya Ramdas and Ruodu Wang.
\newblock Hypothesis testing with e-values.
\newblock arXiv:2410.23614, 2024. DOI: 10.48550/arXiv.2410.23614.
\newblock \url{https://arxiv.org/abs/2410.23614}

\bibitem{wang2022ebh}
Ruodu Wang and Aaditya Ramdas.
\newblock False discovery rate control with e-values.
\newblock \emph{Journal of the Royal Statistical Society: Series B}, 84(3), 822--852, 2022.
\newblock DOI: 10.1111/rssb.12489.

\bibitem{wang2025sebH}
Hongjian Wang, Sanjit Dandapanthula, and Aaditya Ramdas.
\newblock Anytime-valid FDR control with the stopped e-BH procedure.
\newblock arXiv:2502.08539, 2025. DOI: 10.48550/arXiv.2502.08539.
\newblock \url{https://arxiv.org/abs/2502.08539}


% Added in rev0806 citation-closure hygiene pass.
\bibitem{series:state1}
Anonymous.
\newblock State-Dependent Anonymity: Selection and State Evolution as an Attack Surface.
\newblock AnonDHT State series Paper~1, 2026. (This archive.)

\bibitem{series:state1addendum}
Anonymous.
\newblock State-Dependent Anonymity (Addendum): Calibration Vignettes and Time-to-Identification Conversions.
\newblock AnonDHT State series Paper~1A, The-One-True-Archive (2026).

\bibitem{series:state2}
Anonymous.
\newblock Marginally Uniform Committee Contact in Anonymous DHT Lookups: Availability Floors, Perfect-Leakage Counterexamples, and Abort-Safe Equalization.
\newblock Companion preprint in this archive (AnonDHT~State~2).

\end{thebibliography}

\end{document}
